
\SGASubsection{5.9}{Counterexamples}
\SGARefTarget{sga3:xvii:subsection:5:9}

Let us first indicate a procedure for obtaining nontrivial
extensions of a \(k\)-group \(H\) of multiplicative type by a
unipotent \(k\)-group \(U\).  Suppose that an action of \(H\) on
\(U\) is given, and let
\[
  E=U\rtimes H
\]
be the semidirect product.  Let, moreover, a class in
\(\mathrm H^1(k,U)\) be represented by a \(1\)-cocycle \(a\).
The group \(U\) acts on \(E\) by inner automorphisms, and \(a\)
therefore defines a \(k\)-form of \(E\), denoted by \(E_a\).
Suppose in addition that \(U\) is commutative.  Then \(U\) acts
trivially on \(U\) and on the quotient \(E/U=H\), so \(E_a\) is
again an extension of \(H\) by \(U\).

For simplicity, suppose that \(H\) is an étale diagonalizable
group, so that the Galois group
\(\mathscr G=\operatorname{Gal}(\overline{k}/k)\) acts trivially
on \(H(\overline{k})\).  Its action on \(E_a(\overline{k})\) is
then given by
\[
  {}^g(u,h)
  =\bigl({}^g u+a(g)-{}^h a(g),h\bigr)
  \qquad
  \bigl(g\in\mathscr G,\ u\in U(\overline{k}),\
        h\in H(\overline{k})\bigr).
\]
If \(h\in H(k)\), and \(X\) is its inverse image in \(E_a\), then
\(X\) is a torsor under \(U\) defined by the class of the
\(1\)-cocycle
\[
  g\longmapsto a(g)-{}^h a(g)
\]
of \(\mathscr G\) with values in \(U\).  It follows that if there
is a point \(h\in H(k)\) for which this cocycle is nontrivial,
then the extension \(E_a\) is nontrivial.  We apply this
construction in the following two particular cases.

\smallskip
\noindent
\emph{\SGARefLink{sga3:xvii:appendix:A}{(a)} A nontrivial extension of an étale diagonalizable group
\(H\) by \(\mathbb Z/p\mathbb Z\).}
Take
\[
  H=(\mathbb Z/p\mathbb Z)^\times,
  \qquad U=\mathbb Z/p\mathbb Z,
\]
with \(H\) acting on \(U\) by multiplication.  Let \(k\) be a
field of characteristic \(p\), let \(K/k\) be a Galois extension
whose Galois group \(\mathscr G\) is isomorphic to
\(\mathbb Z/p\mathbb Z\), and take a nonzero element
\[
  a\in\mathrm H^1(\mathscr G,U)
   =\operatorname{Hom}(\mathscr G,U).
\]
Then \(E_a\) is the required example.

\smallskip
\noindent
\emph{\SGARefLink{sga3:xvii:section:b}{(b)} A nontrivial extension of an étale diagonalizable group
\(H\) by a smooth connected unipotent group \(U\).}
Choose a nonperfect field \(k\) admitting a \(k\)-form \(U\) of
\(\mathbb G_a\) such that \(\mathrm H^1(k,U)\ne0\).  For example
(see J.-P.\ Serre, \emph{Cohomologie Galoisienne}), one may take
for \(k\) the fraction field of a discrete valuation ring of
equal characteristic \(p>0\), and for \(U\) the algebraic
subgroup of \(\mathbb G_a\times_k\mathbb G_a\) defined by
\[
  X^p+X+tY^p=0,
\]
where \(t\) is a uniformizer of \(k\).  Indeed, suppose for
simplicity that \(k\) contains the \((p-1)\)-st roots of unity.
There is then an exact sequence of algebraic groups
\[
  0\longrightarrow\mathbb Z/p\mathbb Z
   \longrightarrow U
   \longrightarrow\mathbb G_a
   \longrightarrow0,
  \qquad
  (x,y)\longmapsto y,
\]
and hence an exact cohomology sequence
\[
  \mathbb G_a(k)\xrightarrow{\,d\,}
  \mathrm H^1(k,\mathbb Z/p\mathbb Z)
  \longrightarrow\mathrm H^1(k,U)
  \longrightarrow0,
\]
where \(d\) associates to \(x\) the homogeneous space under
\(\mathbb Z/p\mathbb Z\) defined by
\[
  X^p+X+tx^p=0.
\]
Moreover,
\[
  \mathrm H^1(k,\mathbb Z/p\mathbb Z)\simeq k/P(k),
  \qquad P(x)=x^p+x,
\]
and therefore
\[
  \mathrm H^1(k,U)\simeq k/\bigl(P(k)+tk^p\bigr).
\]
If, in addition, \(p\ne2\), then \(t^2\) does not belong to
\(P(k)+tk^p\); consequently \(\mathrm H^1(k,U)\ne0\).

The group \(\boldsymbol\mu_{p-1}\) acts on \(U\) by
\[
  (h,x,y)\longmapsto(hx,hy).
\]
Let \(\mathscr G\) be the Galois group of the extension \(K/k\)
defined by
\[
  X^p+X+t^2=0,
\]
and let \(a\in\mathrm H^1(\mathscr G,U)\) be the nonzero class
described above.  It follows immediately that \(E_a\) is a
nontrivial extension of \(\boldsymbol\mu_{p-1}\) by \(U\).

\smallskip
\noindent
\emph{\SGAInlineRefTarget{sga3:xvii:subsection:5-9:item:c}\SGARefLink{sga3:xvii:appendix:C}{(c)} A nontrivial extension of \(\mathbb G_m\) by
\(\boldsymbol\alpha_p\).}
By \SGARefLink{sga3:xvii:theorem:5:1:1:item:i}{5.1.1(i)}\SGARefLink{sga3:xvii:section:b}{(b)}, such an extension can exist only over a
nonperfect field.  Let \(k\) therefore be nonperfect, and let
\(G\) be the semidirect product of \(U=\mathbb G_a\) by
\(H=\mathbb G_m\), where \(\mathbb G_m\) acts on
\(\mathbb G_a\) by homotheties.  Since \(U\) is invariant in
\(G\), so is \({}_{F}U\).  Put
\[
  G'=G/{}_{F}U.
\]
The group \(G'\) is isomorphic to
\(\mathbb G_a\rtimes\mathbb G_m\), where \(\mathbb G_m\) acts on
\(\mathbb G_a\) by
\[
  (h,u)\longmapsto h^p u.
\]
The functor \(\mathscr T_G\) (respectively
\(\mathscr T_{G'}\)) of subtori of \(G\) (respectively \(G'\))
(see Exposé XV) is isomorphic to \(\mathbb G_a\), and the morphism
\[
  \mathscr T_G\longrightarrow\mathscr T_{G'}
\]
induced by \(G\longrightarrow G'\) identifies with
\[
  u\longmapsto u^p.
\]
It follows that if \(T'\) is a subtorus of \(G'\) corresponding
to a point \(x\in k\simeq\mathbb G_a(k)\) such that
\(x^{1/p}\notin k\), then its inverse image \(E\) in \(G\) is an
extension of the torus \(T'\) by
\({}_{F}U=\boldsymbol\alpha_p\).  It has no maximal torus defined
over \(k\), and hence is nontrivial.  As a subgroup of
\[
  G=\operatorname{Spec}k[U,T,T^{-1}],
\]
the group \(E\) is defined by
\[
  U^p=x-xT^p.
\]

\SGAStatement{Remark 5.9.1}{This last example shows that a
nonsmooth algebraic group defined over a nonperfect field need
not possess a maximal torus, and thus answers the question posed
in \SGARefLink{sga3:xiv:remark:1-5:item:b}{Exposé XIV, 1.5(b)}.}
\SGARefTarget{sga3:xvii:remark:5:9:1}

\smallskip
\noindent
\emph{(d) A trivial extension with nonconjugate lifts.}
We now give an example of a trivial extension \(E\) of a group
\(H\) of multiplicative type by a unipotent group \(U\), and of
two lifts \(H'\) and \(H''\) of \(H\) which are not conjugate by
an element of \(U(k)\).

Take for \(E\) the semidirect product of \(U=\mathbb G_a\) by
\[
  \boldsymbol\mu_p
  =\operatorname{Spec}k[T]/(T^p-1),
\]
where the action of \(\boldsymbol\mu_p\) on
\(\mathbb G_a=\operatorname{Spec}k[U]\) is defined by the
comorphism
\[
  U\longmapsto(U+U^p)T-U^p.
\]
The centralizer of \(\boldsymbol\mu_p\) in \(U\) is then the
étale group \(Z\) defined by
\[
  U+U^p=0.
\]
It follows that if \(k\) is not algebraically closed, the
canonical map
\[
  U(k)\longrightarrow(U/Z)(k)
\]
is not in general surjective.  In view of \SGARefLink{sga3:xvii:theorem:5:1:1:item:ii}{5.1.1(ii)}\SGARefLink{sga3:xvii:section:b}{(b)}, two
lifts of \(\boldsymbol\mu_p\) in \(E\) therefore need not be
conjugate by an element of \(U(k)\).

Here is another example, with \(k\) algebraically closed of
characteristic \(p>0\).  Let \(G\) be the radicial semidirect
product in which \(\boldsymbol\mu_p\) acts on
\(\boldsymbol\alpha_p\) by homotheties.  There is an exact
sequence of algebraic groups with \(\boldsymbol\mu_p\)-operators
\[
  0\longrightarrow\boldsymbol\alpha_p
   \longrightarrow\mathbb G_a
   \xrightarrow{\,x\longmapsto x^p\,}\mathbb G_a
   \longrightarrow0,
\]
where \(\boldsymbol\mu_p\) acts by homotheties on the first
\(\mathbb G_a\) and trivially on the second.  The exact
cohomology sequence (Appendix I, \SGARefLink{sga3:xvii:proposition:a-1-1}{Proposition 11}) gives
\[
  0\longrightarrow\mathbb G_a(k)
   \longrightarrow\mathrm H^1(\boldsymbol\mu_p,\boldsymbol\alpha_p)
   \longrightarrow\mathrm H^1(\boldsymbol\mu_p,\mathbb G_a).
\]
Since the last term is zero by I, \SGARefLink{sga3:i:theorem:5-3-3}{5.3.3}, the group
\(\mathrm H^1(\boldsymbol\mu_p,\boldsymbol\alpha_p)\) is nonzero.
Thus two lifts of \(\boldsymbol\mu_p\) in \(G\) need not be
conjugate.
